% SPDX-License-Identifier: CC-BY-4.0
\documentclass[11pt,letterpaper]{article}

\usepackage[T1]{fontenc}
\usepackage{amsmath,amssymb}
\usepackage{tikz}
\usepackage{pgfplots}
\pgfplotsset{compat=1.18}
\usepackage{algorithm}
\usepackage{algpseudocode}
\usepackage{booktabs}
\usepackage{enumitem}
\usepackage[hidelinks]{hyperref}

\newcommand{\N}{\mathbb{N}}
\newcommand{\Z}{\mathbb{Z}}
\newcommand{\Q}{\mathbb{Q}}
\newcommand{\R}{\mathbb{R}}
\newcommand{\C}{\mathbb{C}}
\newcommand{\Map}[3]{#1 \colon #2 \to #3}
\newcommand{\Mod}[1]{#1\text{-}\mathrm{Mod}}
\newcommand{\GroupHom}{\mathrm{Hom}_{\mathbf{Grp}}}
\newcommand{\RingHom}{\mathrm{Hom}_{\mathbf{Ring}}}

\newtheorem{theorem}{Theorem}
\newtheorem{lemma}[theorem]{Lemma}
\newtheorem{corollary}[theorem]{Corollary}
\newtheorem{definition}[theorem]{Definition}
\newtheorem{remark}[theorem]{Remark}
\newenvironment{proof}{\noindent\textit{Proof.}}{\hfill\(\square\)}

\title{Generalized Collatz Library}
\author{}
\date{\today}

\begin{document}
\maketitle

\begin{abstract}
For integers \(q\ge 2\), \(k\ge 1\), and \(t\ge 1\), the raw map with divisor \(D=q^k\) and offset \(c=tD\) separates positive seeds not divisible by \(q\). Divisibility by \(\gcd(q,t)\) distinguishes an escaping neutral translation
from a positive one-division return bounded by an explicit floor. The escape
bound holds at every raw time. An explicit class-changing return shows why
the floor comparison is not a trapping assertion.
\end{abstract}

\section{The neutral family}

Throughout, \(q\ge 2\), \(k\ge 1\), and \(t\ge 1\) are integers. Write
\[
G_i(q)=1+q+\cdots+q^{i-1}
\]
for \(i\ge 1\), and write \(G_0(q)=0\). Thus \(G_{i+1}(q)=qG_i(q)+1\) and \(G_{i+1}(q)=G_i(q)+q^i\). Set \(D=q^k\), \(c=tD\), and \(b=tG_k(q)\). Let \(g=\gcd(q,t)\) and
\[
A=\frac{q}{g}b-t.
\]
The raw step is \(C(x)=x/D\) when \(D\) divides \(x\), and \(C(x)=qx+c\) otherwise. Iteration is written \(C^0(x)=x\) and \(C^{m+1}(x)=C(C^m(x))\).

\begin{lemma}
\label{lem:shift}
\(b=t+q\cdot tG_{k-1}(q)\), and hence \(b-t=q\cdot tG_{k-1}(q)\). In particular \(b\ge t\ge 1\).
\end{lemma}

\begin{proof}
The identity \(G_k(q)=qG_{k-1}(q)+1\) gives the splitting of \(b\). The second formula is cancellation of \(t\). Since \(G_k(q)\ge G_1(q)=1\), one has \(b\ge t\ge 1\).
Also \(g\mid q\) and \(1\le g\le q\), so \(q/g\) is a positive integer.
It follows that \(A\ge b-t\ge0\).
\end{proof}

\begin{lemma}
\label{lem:same-div}
For every integer \(s\ge 0\), \(q\) divides \(n+sb\) if and only if \(q\) divides \(n+st\). Consequently, if \(g\) does not divide \(n\), then \(q\) divides no integer \(n+sb\).
\end{lemma}

\begin{proof}
Lemma~\ref{lem:shift} gives \(n+sb=n+st+q\cdot stG_{k-1}(q)\). Divisibility by \(q\) is therefore unchanged by the last summand. If \(q\) divides \(n+st\), then \(g\) divides both \(q\) and \(t\), hence divides \(n\). The contrapositive is the second claim.
\end{proof}

\begin{definition}
\label{def:clip}
For a positive integer \(m\), let \(j(m)\) be the largest integer \(j\le k\) such that \(q^j\) divides \(m\). The return of \(m\) is
\[
H(m)=\frac{m}{q^{j(m)}}+tG_{k-j(m)}(q).
\]
\end{definition}

\section{Affine blocks}

\begin{lemma}
\label{lem:affine}
Let \(a\) be a positive integer not divisible by \(q\), and let \(0\le i\le k\). Then
\[
C^i(a)=q^ia+cG_i(q).
\]
If \(i<k\), the right-hand side is not divisible by \(D\). After \(k\) steps it equals \(D(a+b)\), so \(C^{k+1}(a)=a+b\).
\end{lemma}

\begin{proof}
The case \(i=0\) is \(C^0(a)=a\). Suppose the formula holds for some \(i<k\). The second summand of \(q^ia+cG_i(q)\) is divisible by \(D\). If \(D\) also divided the sum, then \(q^{k-i}\) would divide \(a\). Since \(k-i\ge 1\), \(q\) would divide \(a\), a contradiction. Thus the next raw step is affine, and the recurrence for \(G_{i+1}(q)\) gives the formula at \(i+1\). At \(i=k\) the same formula is \(D(a+b)\), so the following raw step divides by \(D\).
\end{proof}

\begin{lemma}
\label{lem:prefix}
Fix \(n\ge 1\) with \(q\nmid n\). Suppose \(q\) does not divide \(n+rb\) for every \(0\le r<s\). Then \(C^{r(k+1)}(n)=n+rb\) for every \(0\le r\le s\).
\end{lemma}

\begin{proof}
At index zero the landing is the seed. If \(r<s\), the anchor \(n+rb\)
is not divisible by \(q\). Lemma~\ref{lem:affine} therefore gives the next
block landing \(n+(r+1)b\). Induction on \(r\) proves the assertion.
\end{proof}

\begin{lemma}
\label{lem:escape}
If \(n\ge 1\), \(q\nmid n\), and \(g\nmid n\), then for every \(s\ge 0\) and every \(i\le k\),
\[
C^{s(k+1)+i}(n)=q^i(n+sb)+cG_i(q).
\]
The raw orbit is unbounded.
\end{lemma}

\begin{proof}
Lemma~\ref{lem:same-div} supplies the hypothesis of Lemma~\ref{lem:prefix} for every \(s\). Lemma~\ref{lem:affine}, applied to the anchor \(n+sb\), gives the displayed formula. Lemma~\ref{lem:shift} gives \(b\ge 1\), so \(C^{(N+1)(k+1)}(n)=n+(N+1)b>N\).
\end{proof}

\begin{corollary}
\label{cor:all-clock}
Under the hypotheses of Lemma~\ref{lem:escape}, for every raw time \(\tau\),
\[
C^\tau(n)\ge n+b\left\lfloor\frac{\tau}{k+1}\right\rfloor.
\]
Consequently \(C^\tau(n)\) tends to infinity as \(\tau\) tends to infinity.
\end{corollary}

\begin{proof}
Euclidean division gives \(\tau=s(k+1)+i\), where \(0\le i\le k\) and
\(s=\lfloor\tau/(k+1)\rfloor\). The formula in Lemma~\ref{lem:escape} is
at least \(n+sb\), because \(q^i\ge1\) and \(cG_i(q)\ge0\). Since \(b\ge1\),
for every \(M\ge0\) and every \(\tau\ge(M+1)(k+1)\), it exceeds \(M\).
\end{proof}

\section{The first divisible anchor}

\begin{lemma}
\label{lem:cover}
Let \(q'>0\) and let \(t'\) be coprime to \(q'\). For every residue \(r\), some \(s<q'\) satisfies \(t's\equiv r\pmod{q'}\).
\end{lemma}

\begin{proof}
Multiplication by \(t'\) on \(\Z/q'\Z\) is injective: a congruence \(t'u\equiv t'v\pmod{q'}\) with \(u\le v\) means that \(q'\) divides \(t'(v-u)\), hence divides \(v-u\), and \(0\le v-u<q'\) forces \(u=v\). An injective self-map of a finite set is surjective.
\end{proof}

\begin{lemma}
\label{lem:delay}
If \(q\nmid n\) and \(g\mid n\), there exists a least integer \(s\) such that \(1\le s\le q/g-1\) and \(q\) divides \(n+sb\).
\end{lemma}

\begin{proof}
Set \(q'=q/g\), \(t'=t/g\), and \(n'=n/g\). Then \(q'>1\): otherwise \(q=g\), so \(q\) divides \(n\). Also \(\gcd(q',t')=1\). Lemma~\ref{lem:cover} supplies \(s_0<q'\) with \(t's_0\equiv -n'\pmod{q'}\), hence \(q'\) divides \(n'+t's_0\). Multiplying by \(g\) shows that \(q\) divides \(n+s_0t\), and Lemma~\ref{lem:same-div} moves the divisibility to \(n+s_0b\). The least such nonnegative integer cannot be \(0\), because \(q\nmid n\), and it is at most \(s_0\), hence at most \(q'-1\).
\end{proof}

\section{The clipped return}

\begin{lemma}
\label{lem:closed}
Let \(m\) be a positive integer and let \(j=j(m)\). For \(0\le i\le k-j\),
\[
C^i(m)=q^im+cG_i(q),
\]
and \(C^{k-j+1}(m)=H(m)\).
\end{lemma}

\begin{proof}
If \(i<k-j\) and \(D\) divided \(q^im+cG_i(q)\), the second summand would force \(q^{k-i}\) to divide \(m\), hence \(q^{j+1}\) to divide \(m\), contradicting maximality of \(j\). The affine recurrence therefore holds through index \(k-j\). At that index the state is
\[
q^{k-j}m+cG_{k-j}(q)=D\cdot H(m),
\]
because \(q^j\) divides \(m\) and \(q^k=q^{k-j}q^j\). The next raw step divides by \(D\). The return is positive: \(q^j\) is
positive and divides the positive integer \(m\), so \(m/q^j\ge1\), while
\(tG_{k-j}(q)\ge0\).
\end{proof}

\begin{lemma}
\label{lem:product}
Suppose \(1\le s\le q/g-1\) and \(q\) divides \(m=n+sb\). Let \(j=j(m)\) and \(J=H(m)\). Then \(Jq\le n+A\).
\end{lemma}

\begin{proof}
Since \(q\mid m\) and \(k\ge1\), exponent one is admissible in
Definition~\ref{def:clip}. Maximality therefore gives \(j\ge1\), so
\(m/q^j\le m/q\). Also \(G_{k-j}(q)\le G_{k-1}(q)\). Therefore
\begin{align*}
Jq
&=q\cdot\frac{m}{q^j}+q\cdot tG_{k-j}(q)
\le m+q\cdot tG_{k-1}(q)
=m+(b-t).
\end{align*}
The bound on \(s\) gives \(m\le n+(q/g-1)b\). Adding \(b-t\) and using \(t\le b\) produces
\[
m+(b-t)\le n+\frac{q}{g}b-t=n+A.
\]
\end{proof}

\begin{lemma}
\label{lem:threshold}
Let \(A,J,n\) be nonnegative integers and suppose \(Jq\le n+A\). Set \(N_0=\lfloor A/(q-1)\rfloor\). If \(N_0<n\), then \(J<n\). If \(n\le N_0\), then \(J\le N_0\).
\end{lemma}

\begin{proof}
The quotient formula gives \(A<n(q-1)\) when \(N_0<n\). Hence \(n+A<nq\). Combined with \(Jq\le n+A\), this yields \(Jq<nq\), so \(J<n\).

If instead \(n\le N_0\), the same quotient formula gives \(A<(N_0+1)(q-1)\), so \(A\le (N_0+1)(q-1)-1\). Then
\[
N_0+A\le N_0+(N_0+1)(q-1)-1=(N_0+1)q-2,
\]
because \((N_0+1)q=(N_0+1)(q-1)+(N_0+1)\). Thus \(n+A<(N_0+1)q\). Combined with \(Jq\le n+A\), this yields \(J<N_0+1\), hence \(J\le N_0\).
\end{proof}

\begin{theorem}
\label{thm:separation}
Let \(n\ge 1\) satisfy \(q\nmid n\). If \(g\nmid n\), the identity and unboundedness of Lemma~\ref{lem:escape} hold. If \(g\mid n\), the delay \(s\) of Lemma~\ref{lem:delay} satisfies \(C^{s(k+1)}(n)=n+sb\), and the return \(J=H(n+sb)\) is the raw state \(k-j(n+sb)+1\) steps later. It obeys \(Jq\le n+A\) and the two comparisons of Lemma~\ref{lem:threshold}.
\end{theorem}

\begin{proof}
The two divisibility conditions on \(g\) are exhaustive and mutually exclusive. The first is Lemma~\ref{lem:escape}. For the second, Lemma~\ref{lem:delay} gives the least delay \(s\), so Lemma~\ref{lem:prefix} identifies the state after \(s\) blocks with \(m=n+sb\). Lemma~\ref{lem:closed} identifies the later return with \(H(m)\). Lemma~\ref{lem:product} gives \(Jq\le n+A\), and Lemma~\ref{lem:threshold} gives the comparisons with \(\lfloor A/(q-1)\rfloor\).
\end{proof}

\section{A small return that escapes}

\begin{theorem}
\label{thm:class-change}
For \(q=4\), \(k=2\) and \(t=2\), the seed \(n=2\) is in the return class
and is below the floor \(N_0=6\). Its return is \(J=5\), and its full raw
orbit tends to infinity. More precisely, for every \(s\ge0\) and \(0\le i\le2\),
\[
C^{5+3s+i}(2)=4^i(5+10s)+32G_i(4).
\]
\end{theorem}

\begin{proof}
Here \(D=16\), \(c=32\), \(g=2\), \(b=10\), and \(A=18\), giving \(N_0=6\).
The first delay is \(s=1\), its landing is \(m=12\), and the clipped exponent
is \(j=1\). The exact raw prefix is
\[
2\longmapsto40\longmapsto192\longmapsto12\longmapsto80\longmapsto5.
\]
The first two and fourth edges are affine; the third and fifth are divisions
by sixteen. Thus \(C^5(2)=5\). Since \(2\nmid5\), the new seed lies in the
escape class. Lemma~\ref{lem:escape} gives its displayed subsequent states,
and Corollary~\ref{cor:all-clock} gives divergence of the full sequence.
\end{proof}

The floor bounds the return in Theorem~\ref{thm:separation}. It does not
assert that the returned seed keeps the original gcd class or that the whole
future orbit is bounded. Theorem~\ref{thm:class-change} separates these claims
with a concrete orbit.

\section{Formal proofs and computation}

The Lean module proves Theorem~\ref{thm:separation}, the all-clock lower bound,
positivity of the clipped return, and the exact class-changing example.
Its source is \texttt{QCD.Discriminator.001/lean/Discriminator.lean}. The residue-cover
argument constructs a finite residue permutation; it does not assume that
\(q\) is prime. This module is part of the default Lean build.

The Go implementation covers
\[
2\le q\le8,\qquad 1\le k\le4,\qquad1\le t\le8,\qquad1\le n\le256,
\qquad q\nmid n.
\]
There are \(43{,}328\) such seeds: \(13{,}440\) escape seeds and \(29{,}888\)
return seeds. Escape checks compare every partial state for block indices
zero through eight with the closed formula. Return checks independently
replay both raw landings and check the product and floor comparisons.
These are finite computations, separate from the unrestricted proofs.

Each supported raw replay uses at most \(44\) ticks; the implementation caps
replay at \(64\). All compared intermediate states are bounded by
\[
4096(256+9\cdot4680)=173{,}572{,}096<2^{28}.
\]
Before using Lemma~\ref{lem:product}, elementary box bounds give
\[
J\le m+tG_k(q)\le37{,}696.
\]
Thus even the tested product \(qJ\) is safely within 64-bit arithmetic. The general arithmetic helpers check overflow; invalid
family parameters are rejected before loop counters or powers are formed.
Constant bases in power and geometric-sum helpers are handled without a
loop. In particular \(G_{64}(2)=2^{64}-1\) is accepted without evaluating the
unneeded power \(2^{64}\).

The tests cover the complete grid, independently locate the first clipped
division, distinguish strict and small returns, and verify the class change.
They also check overflow and unsupported-budget controls. Successful
checks of one seed allocate no memory. Run from the repository root:
\begin{verbatim}
make lean
make go GOFILE=QCD.Discriminator.001/go/Discriminator.go
CGO_ENABLED=0 go test ./QCD.Discriminator.001/go
make tex TEXFILE=QCD.Discriminator.001/tex/Discriminator.tex
bash Scripts/CheckIgnoredPdfs.sh
bash Scripts/CheckIgnoredTexTemps.sh
\end{verbatim}

\end{document}
